\subsection{The same order for ordinary backtracking}
Certificate normalization is not necessary for the two-phase order.
Consider the standard quadratic backtracking scheme of
\cite{pedregosa}, specialized to initial trial $L_{k-1}/2$ and
doubling on rejection, with the same test~\eqref{eq:accept} and
certificate~\eqref{eq:certificate}. Let $L_{-1}=M>0$ be arbitrary.
The algorithm needs neither curvature constants nor a diameter input;
$D$ enters only its analysis. The optional initialization $M=U_0/D^2$
uses a diameter bound and also makes the trajectory invariant under
positive objective rescaling, with the tolerance scaled accordingly.

\begin{theorem}[Ordinary backtracking]\label{thm:ordinary}
Under the assumptions of Theorem~\ref{thm:main}, let
$J=\lceil\log_2(U_0/\varepsilon)\rceil$ and
$P=\lceil\log_2\max\{1,MD^2/U_0\}\rceil$.
Ordinary backtracking reaches $U_N\le\varepsilon$ within
\begin{equation}\label{eq:ordinarycomplexity}
 N\le P+\left\lceil(6+16H_1D^2)J+32H_0D^2/\varepsilon\right\rceil
\end{equation}
updates, using $N+1$ gradients and linear-oracle calls. For the actual
$N\ge1$, all function calls satisfy
\begin{equation}\label{eq:ordinarytrials}
 N_f\le1+2N+
 \left\lceil\log_2\max\{1,2(H_0+H_1U_0)/M\}\right\rceil.
\end{equation}
For $M=U_0/D^2$, $P=0$ and the logarithmic overhead in
\eqref{eq:ordinarytrials} is independent of $\varepsilon$.
\end{theorem}
\begin{proof}
As before, every $L\ge H_0+H_1U_k$ passes the test. If a step has at
least one rejection, its accepted curvature is strictly less than
$2(H_0+H_1U_k)$. If there is no rejection, $L_k=L_{k-1}/2$.
Acceptance implies monotone $f$ and $U$.

Let $b_j=U_0/2^j$, $j=0,\ldots,J-1$, and assign a source iterate to
phase $j$ when $b_j/2<U_k\le b_j$. Sources with $U_k>\varepsilon$
belong to these phases. Define the positive caps
\[
 C_j=\max\{2(H_0+H_1b_j),\ b_j/(2D^2)\}.
\]
They satisfy $C_j/2\le C_{j+1}\le C_j$.
Call a step a lag step if its accepted $L_k>C_j$ in its source phase.
Such a step has no rejection, so its curvature halves. Once
$L_k\le C_j$, all subsequent steps in the same phase also satisfy
this bound: their initial trial is at most $C_j/2$, and a trial
accepted after rejection is below $2(H_0+H_1U_k)\le C_j$.

For completeness, before a step in phase $j$ give the previous estimate
the nonnegative integer potential
\[
 p=\max\{0,\lceil\log_2(L_{k-1}/C_j)\rceil\}.
\]
A lag step reduces this potential by one; a non-lag step leaves
potential zero before any phase change. A jump over $m$ phases can
increase it by at most $m$, because the cap falls by at most $2^m$.
Initially $p\le P+1$, since $C_0\ge U_0/(2D^2)$. Source phase indices
increase by at most $J-1$ in total. Thus at most $P+J$ lag steps occur,
also on any prefix of the run.

On a non-lag step in phase $j$, acceptance, $G_k\ge U_k>b_j/2$,
and $q_k\le D$ give a decrease of at least
\[
 p_j=\min\{b_j/4,\ b_j^2/(8C_jD^2)\}.
\]
All steps decrease $f$, and the total decrease during this phase
is at most $b_j$. Its non-lag steps therefore number at most
\[
 \lfloor b_j/p_j\rfloor+1
 \le 5+16H_1D^2+16H_0D^2/b_j.
\]
Sum over phases and use $\sum_j1/b_j<2/\varepsilon$, then add the
lag bound. The same argument bounds every prefix of an unfinished
run, establishing termination and~\eqref{eq:ordinarycomplexity}.

Finally the accepted estimates obey
\[
 L_k\le\max\{M/2^{k+1},\ 2(H_0+H_1U_0)\}.
\]
Writing $r_k$ for the number of rejections gives
$L_k=2^{r_k-1}L_{k-1}$ and hence exactly
$\sum_{k<N}(r_k+1)=2N+\log_2(L_{N-1}/M)$.
The preceding envelope proves~\eqref{eq:ordinarytrials} after adding
the initial value. Accepted values are cached.
\end{proof}
This is an analysis of an existing algorithm, not a new backtracking
rule. The constants above are conservative and do not imply a practical
ranking of the two variants. A complete literature audit is still needed
to determine whether this specialization or an equivalent argument has
already appeared. In particular, the telescoping trial-count identity
is standard in backtracking analyses~\cite{pedregosa}.

\begin{proposition}[Why an initialization term remains]\label{prop:initialization}
Let $Q=[-1,1]$, $f(x)=-cx$ with $c>0$, $x_0=-1$, and $L_{-1}=M>0$.
Although $H_0=H_1=0$, ordinary halving backtracking needs exactly
\[
 N=\left\lceil\log_2\left(1+\frac Mc
                  \left(1-\frac{\varepsilon}{2c}\right)\right)\right\rceil
\]
updates for $0<\varepsilon<2c$. In particular, halving $U_0$ can require
$\Theta(\log(M/c))$ updates as $M/c\to\infty$.
\end{proposition}
\begin{proof}
Every quadratic test accepts its first trial, since $f$ is affine.
Thus $L_k=M/2^{k+1}$ and
$x_{k+1}=\min\{1,x_k+c/L_k\}$. The linear lower bound is always the
exact optimum $-c$, so
\[
 x_N=\min\{1,-1+(2c/M)(2^N-1)\},\qquad U_N=c(1-x_N).
\]
Solve $U_N\le\varepsilon$ for the integer $N$. This concerns the
initialization penalty of this particular backtracking rule; it is not
a lower bound for arbitrary first-order or linear-oracle algorithms.
With $M=U_0/D^2=c/2$, the first step is full and solves the problem.
\end{proof}
